% ===================================================================
%  TABELLA N-RO 3 — Pueritia e juventute.
%  Traduction 2026 d'apres texto/io, controlee sur texto/fr et
%  texto/es. Consigne : voir texto/ia/10-tabella-01.tex.
% ===================================================================

%%K t03-apar-1 apar
\VUsaut{3.0mm}
\VUtitre{15.0pt}{60}{TABELLA N\textsuperscript{o} 3}
\VUsaut{1.6mm}
\VUtitre{11.4pt}{0}{Pueritia e juventute.}
\VUsaut{3.4mm}

%%K t03-apar-2 apar
(Le 3\textsuperscript{e} e 4\textsuperscript{e} tabellas es assi combinate que
illos presenta in quatro \VUgras{scenas familial} le notiones essential super
le \VUgras{tempore}, le \VUgras{etate}, le \VUgras{familia}, le
\VUgras{vestimento} e le \VUgras{stato}.)
\VUsaut{4.0mm}

%%K t03-tit-1 sub
\VUcentre{12.0pt}{0}{\textit{Prime scena.}}
\VUsaut{3.0mm}

%%K t03-tit-2 sub
\VUpk{10.2pt}{40}{Ceremonia de baptismo in un village.}
\VUsaut{2.4mm}

%%K t03-tit-3 sub
\VUpk{10.2pt}{40}{Le primavera.}
\VUsaut{3.0mm}

%%K t03-c1-01-1 p
1. --- Nos es in le \VUgras{primavera}, in le menses de \VUgras{martio}, de
\VUgras{april} o de \VUgras{maio}, in le dies del \VUgras{septimana},
\VUgras{lunedi}, \VUgras{martedi}, o \VUgras{mercuridi}. Il es \VUgras{dece
horas} del matino.

%%K t03-c1-02-1 p
2. --- Le \VUgras{tempore} es belle, le \VUgras{aere} pur. Le sol brilla. Alcun
\VUgras{parve nubes} \textsuperscript{(2)} ascende in le blau \VUgras{celo}
\textsuperscript{(1)}.

%%K t03-c1-03-1 p
3. --- Le \VUgras{arbores} ha a pena \VUgras{folios}. Le tenue \VUgras{poplos}
\textsuperscript{(3)} e le alte \VUgras{platano} \textsuperscript{(17)} ha
usque ora solmente gemmas.

%%K t03-c1-04-1 p
4. --- Le \VUgras{hirundines} \textsuperscript{(4)} reveni e vola con gai
pipiamento circa le \VUgras{flecha} \textsuperscript{(7)} del \VUgras{campanario}
\textsuperscript{(5)} que corona le \VUgras{ecclesia} \textsuperscript{(6)} del
village.

%%K t03-tit-4 sub
\VUsaut{3.4mm}
\VUpk{10.2pt}{40}{Le ecclesia.}
\VUsaut{3.0mm}

%%K t03-c1-05-1 p
5. --- Multo simple es iste ecclesia con su grande \VUgras{portal}
\textsuperscript{(13)} e grosse \VUgras{horologio} \textsuperscript{(8)}. Le
parve \VUgras{monstrator} \textsuperscript{(9)} es sur le cifra X, illo indica le
\VUgras{horas}. Le \VUgras{grande} \textsuperscript{(10)} es sur XII (meridie);
illo indica le \VUgras{minutas}.

%%K t03-c1-06-1 p
6. --- In le campanario, le \VUgras{campanas} \textsuperscript{(11)} sona
gaimente. Al summitate del tecto sta un parve \VUgras{cruce}
\textsuperscript{(12)} de petra.

%%K t03-c1-07-1 p
7. --- Le ecclesia ha non solmente un grande \VUgras{portal}
\textsuperscript{(13)}, illo ha anque un porta lateral. Tra iste porta le senior
\VUgras{parocho} \textsuperscript{(15)} vestite de su \VUgras{sutana}, exi del
\VUgras{sacristia} \textsuperscript{(14)} e va al \VUgras{presbyterio}
\textsuperscript{(19)}.

%%K t03-c1-08-1 p
8. --- Presso ille, ecce le \VUgras{lactera} \textsuperscript{(24)} con su
\VUgras{asino} \textsuperscript{(23)}. Illa reveni del citate, in le qual illa
portava bon lacte pro le infantes.

%%K t03-tit-5 sub
\VUsaut{4.0mm}
\VUpk{10.2pt}{40}{Le verdiero.}
\VUsaut{3.4mm}

%%K t03-c1-09-1 p
9. --- Senior Aliboron (le asino) es toto contente de iste cursa matinal. Ille
aspira con delicia le aere balsamate del perfumo del \VUgras{flores}
\textsuperscript{(20)} que coperi le \VUgras{arbores fructifere}
\textsuperscript{(18)} del vicin \VUgras{verdiero}. Toto rosate e blanc es iste
\VUgras{persicieros} \textsuperscript{(18)} e \VUgras{abricotieros}
\textsuperscript{(19)}, e illos face sperar un bon recolta.

%%K t03-c1-10-1 p
10. --- Intertanto, le parve \VUgras{apes} \textsuperscript{(22)} del
\VUgras{apiario} \textsuperscript{(21)} va illac pro colliger rapinante lor dulce
\VUgras{melle} susurrante.

%%K t03-c1-11-1 p
11. --- Ma que occurre? Ecce que le pueros del village accurre e face fugir le
\VUgras{ansores} \textsuperscript{(25)} espaventate. Isto es le exito del
ecclesia del cortege post le \VUgras{baptismo} del filio del notario.

%%K t03-c1-12-1 p
12. --- Le parve Jacobo, in su \VUgras{cuna} \textsuperscript{(26)}, se evelia
per le gai sonar del campanas, e su soror Magdalena qui, illa anque, vole vider
le cortege del baptismo, roga su matre venir pro pectinar su capillos e vestir
la al limine del casa.

%%K t03-c1-13-1 p
13. --- Le matre consenti. Illa pone su \VUgras{foulard} \textsuperscript{(28)},
su \VUgras{corsage} \textsuperscript{(29)} e su \VUgras{avantal}
\textsuperscript{(30)}, e illa veni sider sur un parve sedia ante le porta.
Magdalena ha solmente su \VUgras{subgonna} \textsuperscript{(31)} e su
\VUgras{corsetto} \textsuperscript{(32)}.

%%K t03-c1-14-1 p
14. --- Quanto felice illa es! Illa oblida su flores favorite, su
\VUgras{dianthos} \textsuperscript{(34)}, sur le quales se pone le
\VUgras{papiliones} \textsuperscript{(35)}, su \VUgras{tulipas}
\textsuperscript{(36)}, su \VUgras{convallarias} \textsuperscript{(37)}, in
\VUgras{potes} \textsuperscript{(38)}, sur le \VUgras{muro}
\textsuperscript{(33)}, e su \VUgras{rosas} \textsuperscript{(39)} que forma un
haga tanto belle. Illa contempla le ric vestimentos del damas del escorta.

%%K t03-c1-15-1 p
15. --- Le pueros del village non presta multe attention al vestimentos. Illes
se precipita e se pulsa le unes le alteres pro haber \VUgras{bonbones}
\textsuperscript{(55)} o \VUgras{monetas} \textsuperscript{(52)}. Un de illes
plora \textsuperscript{(40)}.

%%K t03-c1-16-1 p
16. --- Le altere calcava sur le pede del \VUgras{can} \textsuperscript{(42)}
qui criante fugi e face fugir le \VUgras{gallina} \textsuperscript{(43)} e su
\VUgras{pullos} \textsuperscript{(44)}.

%%K t03-tit-6 sub
\VUsaut{5.0mm}
\VUpk{10.2pt}{40}{Le cortege del baptismo.}
\VUsaut{4.0mm}

%%K t03-c1-17-1 p
17. --- Ante le escorta marcha le \VUgras{nutrice} \textsuperscript{(45)}. Illa
ha un belle \VUgras{coiffa} \textsuperscript{(46)} con longe bandas. Illa porta
le \VUgras{neonato} \textsuperscript{(47)} toto vestite de \VUgras{punctos}
\textsuperscript{(48)}.

%%K t03-c1-18-1 p
18. --- Detra le nutrice veni le \VUgras{patrino} \textsuperscript{(49)} e le
\VUgras{matrina} \textsuperscript{(53)}. Illes distribue le bonbones e monetas.
Le patrino ha \VUgras{guantos} \textsuperscript{(51)} e un \VUgras{cappello
cylindric} \textsuperscript{(50)}. Le matrina tene in le mano un belle
\VUgras{bouquet} \textsuperscript{(54)}.

%%K t03-c1-19-1 p
19. --- Detra illes, nos vide le \VUgras{oncle} \textsuperscript{(56)} e le
\VUgras{tante} \textsuperscript{(58)} del infante. Le tante ha un elegante
\VUgras{cappello} \textsuperscript{(59)}, le oncle un nigre \VUgras{redingote}
\textsuperscript{(57)} e un blanc gilet. Ecce postea le \VUgras{cosino}
\textsuperscript{(60)} e le \VUgras{cosina} \textsuperscript{(61)}. Iste ultime
tene per le mano le \VUgras{soror} \textsuperscript{(62)} del neonato, le parve
Berta.

%%K t03-c1-20-1 p
20. --- Le altere \VUgras{fratre} \textsuperscript{(63)} de Berta marcha detra.
Ille da le mano a un \VUgras{parenta} \textsuperscript{(64)}. Le \VUgras{amicos}
\textsuperscript{(65)} del familia veni postea.

%%K t03-tit-7 sub
\VUsaut{4.6mm}
\VUpk{10.2pt}{40}{Le spectatores.}
\VUsaut{3.4mm}

%%K t03-c1-21-1 p
21. --- Presso le cortege, ecce un \VUgras{mendicante}
\textsuperscript{(66)} appoiate sur su \VUgras{cruciata} \textsuperscript{(67)}.
Ille demanda un almosna.

%%K t03-c1-22-1 p
22. --- Al altere latere il ha un parve puero multo \VUgras{polite}
\textsuperscript{(68)}. Ille saluta e desira « \VUgras{bon die} » al personas
que ille cognosce.

%%K t03-c1-23-1 p
23. --- Le villageses exiva de lor casas pro vider le cortege del baptismo. Nos
vide un \VUgras{campanio} \textsuperscript{(69)} con su \VUgras{bonetto de
coton} e presso ille un \VUgras{campania} \textsuperscript{(70)}. Postea un
\VUgras{ancian} \textsuperscript{(71)} appoiate sur su \VUgras{baston}
\textsuperscript{(72)}. Finalmente le \VUgras{molinero} \textsuperscript{(73)}
con su grande \VUgras{cappello de feltro} \textsuperscript{(74)} e un juvene
campania calceate de \VUgras{sabots de ligno} \textsuperscript{(75)}, qui
insenia marchar a su infantetto.

%%K t03-c1-24-1 p
24. --- Iste scena es multo amusante.
\VUsaut{4.0mm}

%%K t03-tit-8 sub
\VUcentre{12.0pt}{0}{\textit{Secunde scena.}}
\VUsaut{3.0mm}

%%K t03-tit-9 sub
\VUpk{10.2pt}{40}{Festa public.}
\VUsaut{2.4mm}

%%K t03-tit-10 sub
\VUpk{10.2pt}{40}{Jocos nautic e regatas.}
\VUsaut{3.0mm}

%%K t03-c2-01-1 p
1. --- Le \VUgras{estate} arrivava. Le multo calide sol del mense de
\VUgras{junio} (julio o \VUgras{augusto}) incita le juvenes a banar se e
canotar.

%%K t03-c2-02-1 p
2. --- Sur le ripa dextre del fluvio, le canoteros se disvesti pro participar in
le jocos nautic e regatas.

%%K t03-c2-03-1 p
3. --- Un de illes leva su \VUgras{scarpas} \textsuperscript{(1)}; ille los
disliga (disbottona). Su duo camerados ja es preste. De ora illes ha solmente
lor \VUgras{maillot} \textsuperscript{(2)} e \VUgras{calzon de banio}
\textsuperscript{(3)}. Illes conversa attendente lor amicos. Illes parla del
feria e del festa que ha loco sur le altere ripa.

%%K t03-c2-04-1 p
4. --- Sur le solo, presso illes, jace le \VUgras{vestimentos} de lor camerados
: un \VUgras{par} de \VUgras{bottinas} \textsuperscript{(4)}, \VUgras{scarpas}
\textsuperscript{(5)}, \VUgras{calcettas} \textsuperscript{(6)}, cappellos de
\VUgras{feltro} \textsuperscript{(7)} e de \VUgras{palea} \textsuperscript{(8)}
e un \VUgras{cravata} \textsuperscript{(9)}.

%%K t03-c2-05-1 p
5. --- Un del juvenes plicava curosemente su \VUgras{jachetta}
\textsuperscript{(10)}. Ille lo pone sur un toalia, in le qual su vicino tosto
inserrara le duo \VUgras{vestimentos}.

%%K t03-c2-06-1 p
6. --- Un sexte puero es in retardo. Ille ha ancora su \VUgras{casquetta}
\textsuperscript{(11)} sur le capite. Ille non levava ni su \VUgras{camisa}
\textsuperscript{(12)}, ni su \VUgras{collar} \textsuperscript{(13)}, ni su
\VUgras{manchettas} \textsuperscript{(14)}. Ille tene in le mano su
\VUgras{gilet} \textsuperscript{(17)} e ha ancora su \VUgras{pantalon}
\textsuperscript{(16)} e \VUgras{bracieras} \textsuperscript{(15)}. Certemente
ille non essera preste pro le cursa proxime.

%%K t03-c2-07-1 p
7. --- Presso le juvenes, un parve via al lateres del qual il ha multe
\VUgras{cardos} \textsuperscript{(18)}, conduce al bordo del riviera, ubi le
patre Jacobo prepara inter le \VUgras{cannas} \textsuperscript{(19)} le
\VUgras{artificios} \textsuperscript{(20)} que on tirara le vespere.

%%K t03-tit-11 sub
\VUsaut{4.4mm}
\VUpk{10.2pt}{40}{Le cursa.}
\VUsaut{3.4mm}

%%K t03-c2-08-1 p
8. --- Sur le riviera, alcun damas, protegente se con lor parasoles, promena in
un \VUgras{barca} \textsuperscript{(21)}. Illas seque le cursa que es multo
interessante. In cata \VUgras{yole} \textsuperscript{(22)}, quatro
\VUgras{remeros} \textsuperscript{(24)} rema con tote lor fortias, durante que
le \VUgras{timonero} \textsuperscript{(26)} tene le \VUgras{timon}
\textsuperscript{(25)} e dirige le fragile barca a remos. Le \VUgras{remos}
\textsuperscript{(23)} batte le aqua in cadentia e le legier barcas naviga con
un rapiditate vertiginose.

%%K t03-c2-09-1 p
9. --- Alcun juvenes lucta inter se, presso le pilastro del ponte, pro le
premio del \VUgras{mast unctate} \textsuperscript{(28)}. Un de illes marcha con
prudentia sur le flexibile e lubric mast pro attinger le parve \VUgras{bandiera}
\textsuperscript{(29)} fixate al extremitate. Su infortunate
\VUgras{concurrente} \textsuperscript{(30)} cadeva in le aqua. Ille nata pro
attinger le terra.

%%K t03-c2-10-1 p
10. --- Juvenes plus etose face un \VUgras{torneo (de barcas)}
\textsuperscript{(32)} presso le \VUgras{quai} \textsuperscript{(33)}. A tote
iste jocos assiste un numerose \VUgras{publico} \textsuperscript{(31)} appoiate
al \VUgras{parapetto} del \VUgras{ponte} \textsuperscript{(27)} e amassate sur
le \VUgras{ripa}. Alcun spectatores, tamen, se torna pro sequer le joco del
\VUgras{mast de premio} \textsuperscript{(45)} o pro admirar le \VUgras{escorta
municipal} que veni pro le \VUgras{distribution} del premios.

%%K t03-c2-11-1 p
11. --- Primo, ecce alcun \VUgras{garsonettos} \textsuperscript{(35)} qui
preceede le \VUgras{musicos} \textsuperscript{(36)}; postea veni le \VUgras{maire}
\textsuperscript{(37)}, le adjunctos e le \VUgras{consilio municipal} del parve
citate. Finalmente marcha le \VUgras{pumperos} \textsuperscript{(38)}.

%%K t03-tit-12 sub
\VUsaut{5.0mm}
\VUpk{10.2pt}{40}{Le feria.}
\VUsaut{3.6mm}

%%K t03-c2-12-1 p
12. --- Numerose personas es sur le \VUgras{placia del feria}
\textsuperscript{(39)}. On veni vider le diverse \VUgras{barracas} : le
\VUgras{cavallos de ligno} \textsuperscript{(40)}, le \VUgras{circo}
\textsuperscript{(41)} ubi le representation ja comenciava; le \VUgras{ballo
public} \textsuperscript{(42)}, ubi le juvenes e le juvenas del citate gaude le
placer del \VUgras{dansa}.

%%K t03-c2-13-1 p
13. --- In le fundo, nos vide le \VUgras{arena} \textsuperscript{(44)} ubi le
brave \VUgras{matador} espaniol lucta contra le furiose \VUgras{tauro}
\textsuperscript{(45)}, postea le \VUgras{pista de patinage}
\textsuperscript{(49)}, le \VUgras{stand de tiro} \textsuperscript{(46)} ubi on
se exerce a usar le \VUgras{armas de foco} e a tirar al \VUgras{scopo}.

%%K t03-c2-14-1 p
14. --- Presso, ecce le \VUgras{theatro de feria} \textsuperscript{(47)} in le
qual on joca le \VUgras{comedia} e le \VUgras{drama}. Multo presso, il ha le
barraca del \VUgras{gigante} \textsuperscript{(48)} e del \VUgras{nano}. Le
\VUgras{statura} del prime es duo metros plus dece-cinque centimetros; le
secunde es a pena tanto grande como un infante.

%%K t03-c2-15-1 p
15. --- Finalmente, ecce le grande \VUgras{menagerie} \textsuperscript{(50)} con
su \VUgras{bestias salvage}, su \VUgras{tigre} \textsuperscript{(51)} con
potente \VUgras{ungues}, su \VUgras{leon} \textsuperscript{(52)} con dense
\VUgras{criniera}, su duo \VUgras{girafas} \textsuperscript{(53)} e su
\VUgras{elephante} \textsuperscript{(54)} qui usa tanto habilemente su
\VUgras{trompa}.

%%K t03-c2-16-1 p
16. --- Le \VUgras{domator} \textsuperscript{(55)} qui dressa iste bestias es al
porta del barraca.
\VUsaut{6.0mm}
\VUfilet{22mm}
